% !TeX root = main.tex
\section{长期决策准则的序贯结构和求解条件}\label{sec:structure}

长期决策准则规定了完整策略之间的比较方式，从它出发进行序贯求解还要满足一些结构条件。首先要判断原准则能否等价分解为逐阶段可加的标量目标；如果不能，就要确定决策时必须保留哪些历史信息、价值函数如何表示未来结果；在此基础上，再考察这种表示是否支持递推优化，以及最优性结论依赖哪一类策略。本章依次讨论这四个问题。

\subsection{长期准则的分解与等价转化}\label{sec:structure-decomp}

固定线性权重与期望可加回报是可以直接转化的典型情形。设各目标采用相同的折扣因子，偏好由固定权重向量 \(\mathbf w\) 表示，则在相关期望存在时有
\begin{equation}
\mathbf w^\top
\mathbb E_\pi\!\left[
\sum_{t=0}^{T-1}\gamma^t\mathbf r_t
\right]
=
\mathbb E_\pi\!\left[
\sum_{t=0}^{T-1}\gamma^t
\mathbf w^\top\mathbf r_t
\right].
\end{equation}
定义标量即时回报 $r_t^{\mathbf w}=\mathbf w^\top\mathbf r_t$，原目标就变成期望标量累计回报。若环境状态满足 Markov 条件，且问题中没有需要单独保留的约束或比较规则，就可以采用单目标规划或学习方法\cite{hayes2022}。

某些形式上非线性的 SER 效用也属于这一情形。若
\begin{equation}
u(\mathbf x)=h(\mathbf w^\top\mathbf x+b),
\end{equation}
且 $h$ 严格递增，则它与仿射函数 $\mathbf w^\top\mathbf x+b$ 给出相同的策略排序，仍可优化固定权重下的期望累计回报\cite{rodriguezsoto2026utility}。这里保持的是 SER 排序，不需要把 $h$ 分配到逐阶段回报上。ESR 则不同：$\mathbb E[h(\mathbf w^\top\mathbf G^\pi+b)]$ 一般不由 $\mathbb E[\mathbf w^\top\mathbf G^\pi]$ 决定，这一等价关系在 ESR 下不成立。

可见，能否转化取决于目标函数的分解结构，而不只取决于最终评价结果是否为标量。对于作用于累计结果的非线性效用，一般有
\begin{equation}
u_p\!\left(
\sum_t\gamma^t\mathbf r_t
\right)
\neq
\sum_t\gamma^t u_p(\mathbf r_t).
\end{equation}
把累计结果上的非线性效用替换为逐阶段效用之和，会改变不同阶段结果之间的补偿关系，进而可能改变策略排序。反过来，如果问题本身规定的就是各阶段效用的折扣累积，阶段效用可以直接作为标量奖励。因此，判断能否转化为标准可加目标时，先区分效用作用于阶段结果还是累计结果，再看具体函数的结构。

优先级与约束在转化中也要保持原有语义。严格词典序关系一般无法由有限的大权重代替，硬约束也不同于奖励函数中的惩罚项\cite{skalse2022,achiam2017}。即使主目标已经写成标量累计回报，只要约束或优先级仍独立参与策略比较，问题就仍然是带约束或分层的决策问题，而不是标准的单目标问题。

\subsection{状态充分性与价值表示}\label{sec:structure-state}

环境状态描述了系统未来的演化规律，但不一定包含长期评价所需的全部信息。在标准可加准则下，已经获得的回报不影响未来动作之间的比较；当效用作用于全程累计结果时，历史结果会直接改变当前动作的相对价值。两个执行过程即使到达相同的环境状态，只要此前在各目标上的累计结果不同，后续动作的偏好排序就可能不同。

这时可以把与长期评价有关的历史量纳入状态。设
\begin{equation}
\mathbf c_t
=
\sum_{k=0}^{t-1}\gamma^k\mathbf r_k,
\end{equation}
定义扩展状态
\begin{equation}\label{eq:augmented-state}
\tilde s_t=(s_t,\mathbf c_t,t),
\end{equation}
其中 \(\mathbf c_t\) 是截至时刻 \(t\) 已实现的累计结果，时间变量记录剩余时域及折扣位置。对于路径预算、资源消耗或轨迹事件等准则，还可以加入剩余预算、事件是否已发生等变量。扩展状态需要保留什么由具体准则决定，并不要求保存完整历史\cite{bauerle2014,hayes2023}。


状态扩展恢复了评价所需的信息，也提高了求解复杂度。即使原始状态空间有限，累计回报、剩余资源等变量也可能构成连续或高维空间，增加规划、函数逼近和数据覆盖的难度。因此，状态设计要找的不是尽可能完整的历史，而是足以支持原准则评价与决策的最小表示。

状态决定决策时可以利用哪些信息，价值表示则决定如何总结未来结果。对于 SER，给定策略的期望回报向量就足以计算效用；对于一般非线性 ESR，只保留期望回报向量通常不足以确定期望效用，可能需要表示累计回报分布、与效用有关的统计量，或者直接估计扩展状态上的条件期望效用\cite{hayes2023,bauerle2014}。

不过，非线性准则并不一定要求显式学习完整分布。所需的价值表示取决于评价函数依赖哪些结果信息：有些准则只依赖有限个统计量，有些准则可以用扩展状态上的标量价值函数计算。判断标准是该表示能否支持原准则下的正确比较。多个候选策略之间如何依据回报分布保留和剪枝，在第~\ref{sec:solving-reuse}~节讨论。

\subsection{递推关系与动态一致性}\label{sec:structure-recursion}

有了充分的信息表示，还要判断长期准则是否支持递推求解。这里有两个层次：给定策略的价值能否递推评价，以及最优策略能否通过逐阶段的局部优化获得。前者是策略评价问题，后者对应 Bellman 最优性关系，两者并不等价。

标准期望可加回报下，两类递推都有成熟的形式。对于更一般的准则，阶段分解、条件期望与优化运算之间的关系需要重新考察。例如，在非线性 SER 下，给定策略的期望向量回报可以逐分量递推计算，但不能在每个状态上对向量价值施加效用函数后独立选择动作：局部最优选择是否与初始状态下的整体目标一致，取决于具体准则的结构\cite{roijers2013}。

对于作用于累计结果的非线性 ESR，原状态空间中的标准可加形式不再成立，但在扩展状态 \eqref{eq:augmented-state} 上可以重新构造递推。Bäuerle 和 Rieder 研究标量累计成本 $C^\pi$ 的确定性等价 $U^{-1}(\mathbb E_\pi[U(C^\pi)])$；由于 $U^{-1}$ 严格递增，最小化时可以去掉外层逆函数，只保留累计成本的期望效用\cite{bauerle2014}。在连续严格递增效用、有界正成本、紧的可行动作集合及相应半连续条件、弱连续转移核等假设下，无折扣有限时域问题可以用环境状态与累计成本 $(s,y)$ 构造动态规划，以 $U(y)$ 为终端值，并取得最优决策规则。有折扣时还要跟踪折扣位置：若阶段成本为 $\ell(s,a)$、折扣为 $\beta$，一般效用下扩展变量满足 $y'=y+\zeta\ell(s,a)$、$\zeta'=\beta\zeta$。特殊效用可以简化状态，例如无折扣有限时域的指数效用中，价值里的累计成本因子可以分离并从动作比较中消去；这一简化依赖函数结构和无折扣设定。

把同样的构造用于多目标有限时域 ESR。若状态转移和回报的条件分布由当前状态与动作决定，在扩展状态上定义终端价值
\begin{equation}
V_T(s,\mathbf c)=u_p(\mathbf c),
\end{equation}
并在期望有限、最优动作可取得等条件下有
\begin{equation}
V_t(s,\mathbf c)
=
\max_a
\mathbb E\!\left[
V_{t+1}
\left(
S_{t+1},
\mathbf c+\gamma^t\mathbf R_t
\right)
\mid
S_t=s,A_t=a
\right].
\end{equation}
非线性效用始终作用在完整累计结果上，没有被拆成阶段奖励。部分在原状态空间中不具备可加形式的问题，由此通过状态扩展恢复了递推结构，代价是更大的状态空间和更复杂的价值表示。学习方法面临同样的信息需求：Reymond 等在策略梯度推导中把已累计回报与未来回报一起代入效用，并借助扩展状态形成条件价值表示\cite{reymond2023}。

另一条路线是直接在轨迹分布上规定比较关系。Colaço Carr 等在有限动作、有限观测和有限时域下证明，若可达轨迹分布上的偏好为全预序，并满足“加入同一分布的混合后弱偏好不逆转”的一致性条件，则存在确定性的历史依赖最优策略，且各可达历史下的递归比较给出了最优性的充分条件\cite{carr2024}。这一条件不要求偏好有阶段可加的标量奖励表示，一些不满足期望效用表示条件的偏好也有最优策略。因此，递推结构能否存在，更一般地取决于长期偏好与子问题之间的比较关系是否一致，而不只取决于奖励能否写成显式加和。

可递推之外还有动态一致性问题：初始时刻制定的后续决策，在获得新信息、重新评价后是否仍符合原准则。如果重新决策时忽略已实现的累计结果，或者把初始时刻对完整轨迹的评价换成当前时刻同形式的条件评价，原问题的决策语义就会改变。

风险敏感准则中这一问题尤为突出。对完整结果做一次性风险评价，与逐阶段嵌套的条件风险评价，一般不是同一个目标。Ruszczyński 将时间一致性与条件风险算子的嵌套结构联系起来，并据此构造风险厌恶动态规划\cite{ruszczynski2010}。可递推性回答的是长期准则能否分解求解，动态一致性则要求这种分解在重新决策时仍保持跨期偏好的含义。状态扩展补充的是缺失的信息，它本身并不保证任意评价算子的动态一致性。

\subsection{策略形式与最优性条件}\label{sec:structure-policy}

长期准则还会影响实现最优解所需的策略类别。有限状态、有限动作、奖励有界的无限时域折扣 MDP 存在确定性平稳最优策略，这一结论建立在标量可加准则之上，用于非线性评价、有限时域或带额外约束的多目标问题时需要重新论证\cite{roijers2013}。是否允许随机策略、非平稳策略和历史依赖策略，因此是问题定义的一部分\cite{hayes2022}。

对于有限状态与动作、回报有界且折扣因子小于 $1$ 的 MOMDP，连续 SER 效用给出了一类存在性条件：固定起点后，可达期望回报向量构成紧集，连续效用在其上取得最大值，并可由平稳随机策略实现\cite{rodriguezsoto2026utility}。实现策略可能是随机的，而且原状态上的平稳映射不一定对所有起点同时最优。连续性是充分而非必要条件，但完全去掉它，最大值可能不存在。例如，一阶段的两个动作分别给出 $(1,0)$ 和 $(0,1)$，以前者的选择概率 $p$ 表示随机策略，取
\begin{equation}
u(x_1,x_2)=\frac{x_1}{2}+x_2+\mathbf 1_{\{x_1>0\}}.
\end{equation}
则 $p>0$ 时 SER 为 $2-p/2$，$p=0$ 时为 $1$；效用的上确界为 $2$，但没有策略能达到\cite{rodriguezsoto2026utility}。定义了结果排序，并不保证相应的策略优化问题有最优解。

随机化的作用同样取决于准则。第~\ref{sec:criteria-stochastic}~节例子中的策略 A，可以看作在每次执行开始时等概率选择两个确定性策略，二者分别固定产生 $(2,0)$ 和 $(0,2)$。在 $u=\min\{x_1,x_2\}$ 下，两个确定性策略的 SER 都是 $0$，混合后的期望回报为 $(1,1)$，SER 升至 $1$；而在 ESR 下，每次执行仍只得到 $(2,0)$ 或 $(0,2)$，混合策略的期望效用仍为 $0$。因此，策略混合在 SER 下可以优于任一组成策略，在 ESR 下则没有这种作用。此外，在执行开始时随机选择完整策略，与在执行过程中逐状态随机选择动作，能实现的结果集合也不一定相同\cite{roijers2013}。

历史依赖、Markov 性和平稳性需要分别界定。若策略在同一时刻、相同环境状态下，因此前累计结果不同而采取不同动作，那么相对于原始环境状态，它是历史依赖策略；若累计回报等历史信息可以用扩展变量充分表示，同一决策规则又可以写成扩展状态上的 Markov 策略。平稳性进一步要求策略映射不显式依赖时间。有限时域的最优策略即使是 Markov 的，通常也随剩余时域变化；把时间放进扩展状态后，策略可以统一写成扩展状态到动作的映射，但在原始状态空间上它仍是非平稳的\cite{bauerle2014}。讨论策略类别时，要同时说明状态变量包含哪些历史信息、时间依赖如何表示。

严格词典序在一定条件下仍保留确定性平稳最优策略。Shakerinava 等用偏好公理刻画了这类条件：除完备性、传递性与独立性外，还要求共同的事件前缀不改变后续轨迹彩票的排序，或者该事件使后续结果完全无关。相应的词典序效用可以表示为
\begin{equation}
\mathbf U(e\cdot\tau)=\mathbf r(e)+\Gamma(e)\mathbf U(\tau),
\end{equation}
其中 $e$ 为事件，$\tau$ 为后续轨迹，$\Gamma(e)$ 为对角元为正的下三角矩阵，终止事件对应零矩阵\cite{shakerinava2025lex}。在有限词典序 MDP 中，若所有事件的对角元都小于 $1$，则存在对所有起点同时最优的确定性平稳策略。证明是逐级优化的：先保留最高优先级目标的最优动作，再在受限动作集合中优化下一目标；此时下三角结构中较高优先级的价值已经固定，可以并入下一层回报。严格优先关系由动作集合的逐层限制来保持。该折扣条件有实际边界，例如论文中安全优先的专门模型第一个对角元为 $1$，不满足这一存在性定理的条件。

实际部署条件也会限制允许的策略类别，例如系统不允许随机化，或只能保存有限记忆。这些要求应直接写进策略集合的定义，最优性结论也应说明针对哪一类策略成立。多目标长期平均问题是一个具体例子：Brázdil 等证明，在一般有限 MDP 中，实现给定的期望平均回报向量可能同时需要记忆与随机化\cite{brazdil2014}，标准折扣 MDP 的策略存在性结论在这里不适用。

路径事件约束带来另一种记忆需求。Etessami 等研究让同一策略使多个 $\omega$ 正则或线性时序逻辑性质的满足概率分别达到给定阈值的问题\cite{etessami2008modelchecking}。算法先把时序性质转为乘积模型上的多目标可达性，再用线性规划判断阈值向量是否可达并构造策略。在目标集合均为吸收态的可达性模型中，随机无记忆策略已经足够；映射回原时序性质时，乘积状态记录的信息表现为有限记忆，随机化与记忆都可能是必需的。这里约束的是各路径性质分别成立的概率，与期望累计回报约束、或要求多个事件在同一轨迹上同时发生，语义都不相同。

\subsection{从结构条件到求解方法}\label{sec:structure-to-solving}

本章的结构条件决定了求解方法必须保留什么：能等价转化为期望可加标量回报的准则，可以直接使用单目标方法；依赖累计结果、轨迹特征或结果分布的准则，需要状态扩展或分布表示；优先级、约束和特定策略类别属于问题定义时，必须在优化过程中一直保留。具体采用哪条求解路线，还取决于模型是否可用、能否新增交互、历史数据覆盖了哪些状态与动作。下一章在这三类信息条件下讨论具体方法，并在第~\ref{sec:solving-summary}~节总结准则、信息条件与方法之间的对应。
